\setcounter{neje}{24}
\begin{nejer}
Muestre que el conjunto de Cantor en $\mathbb{R}$ es un conjunto Jordan-nulo. Muestre que $\mathbb{Q} \cap[0,1]$ es un conjunto Lebesgue-nulo, pero no un conjunto Jordan-nulo. Muestre que $A \subseteq \mathbb{R}^n$ es un conjunto Lebesgue-nulo si y solo si $\nu(A) = 0$. Encuentre un conjunto abierto $U \subseteq \mathbb{R}$ cuya frontera tenga medida de Lebesgue positiva.
\end{nejer}
\begin{sol}
% Sean $0<\alpha<1$ y $r\coloneq\frac{1-\alpha}{2}$. En cada paso, se elimina una proporción $\alpha$ del centro de cada intervalo. Este proceso se hace mediante la siguiente recursión:
% \[
% C^{\alpha}_0=[0,1]
% \quad\text{y}\quad 
% \forall n\en\mathbb{N},\; C^{\alpha}_{n+1}=rC^{\alpha}_n \dotcup (1-r + rC^{\alpha}_n)
% \]
% El \emph{$\alpha$-conjunto de Cantor} se define: $\mathcal{C}_\alpha = \bigcap_{n=0}^{\infty} C^{\alpha}_n$, y el \emph{conjunto de cantor} es: $\mathcal{C} \coloneq \mathcal{C}_{1/3}$.\par 
% Como las funciones $a \mapsto  ra$ y $a \mapsto 1-r + a$ son homeomorfismos y $C^{\alpha}_0$ es compacto, entonces también lo es cada $C^{\alpha}_n$, de donde se sigue que $\mathcal{C}_\alpha$ es compacto al ser intersección de compactos; además, aplicando inducción se sigue de inmediato que cada $C^{\alpha}_{n+1}$ es unión de $2^n$ intervalos disjuntos de longitud $r^n$.
\emph{Repasemos la definición de un $\alpha$-conjunto de Smith--Volterra--Cantor}:\\
Sea $0 < \alpha \leq \nicefrac{1}{3}$. Definimos inductivamente una sucesión de
conjuntos compactos $\set{C^\alpha_n}_{n\ge 0}$ en $[0,1]$ del siguiente modo:
\begin{itemize}
\item $C^\alpha_0 = [0,1]$.
\item Suponiendo que $C^\alpha_{n-1}$ es unión de $2^{n-1}$ intervalos cerrados
disjuntos con la misma longitud, $C^\alpha_{n}$ se obtiene eliminando de cada uno de ellos el sub-intervalo abierto
central de longitud $\alpha^{n}$.
\end{itemize}
El \emph{$\alpha$-conjunto de Smith--Volterra--Cantor} se define como $\mathcal{C}_\alpha \coloneq \bigcap_{n=0}^{\infty} C^\alpha_n$, y el \emph{conjunto de Cantor} como $\mathcal{C} \coloneq \mathcal{C}_{1/3}$. Se sigue de inmediato que $\mathcal{C}_\alpha$ es compacto (pues es intersección de cerrados dentro de un compacto). Más aún, si denotamos $S_n$ a la suma de las longitudes de todos los intervalos en $C^\alpha_n$, tenemos que $S_{n}= S_{n-1} - 2^{n-1}\alpha^{n}$ (ya que para obtener $C^\alpha_{n}$ se restaron $2^{n-1}$ intervalos de longitud $\alpha^{n}$), así 
\[
\forall m\en\mathbb{N},\quad
S_{m} - S_0 = \sum^m_{n=1} S_{n} - S_{n-1} = - \sum^m_{n=1}2^{n-1}\alpha^{n}
\]
Como $S_0 = 1$ y $\sum^m_{n=1}2^{n-1}\alpha^{n} = \alpha\sum^m_{n=1}(2\alpha)^{n-1} = \alpha\frac{(2\alpha)^m - 1}{2\alpha-1}$ (pues $2\alpha \le \nicefrac{2}{3} < 1$), se sigue que 
\[
S_m = 1 - \alpha\frac{(2\alpha)^m - 1}{2\alpha-1}
\]
Además, $C^\alpha_m$ es unión de $2^m$ intervalos de la misma longitud, por lo tanto, la longitud de cada intervalo será $2^{-m}S_m$.\vspace{1em}\\
\emph{Probemos que $\mathcal{C}$ es Jordan-nulo}:\\
Sea $\varepsilon>0$. Como $0 < \nicefrac{2}{3} < 1$, existe $N\en\mathbb{N}$ tal que $(\nicefrac{2}{3})^N < \varepsilon$. Luego, ya vimos que $C^{1/3}_N = \dot\bigcup^{2^N}_{k=1} I_k$, donde los $I_k$ son intervalos cerrados disjuntos cuya suma de longitudes (volúmenes) es 
\[
S_N = 1 - \frac{1}{3}\cdot\frac{(\frac{2}{3})^N - 1}{\frac{2}{3}-1} = \left(\frac{2}{3}\right)^N
\]
Por lo tanto,
\[
\mathcal{C} \subseteq C^{1/3}_N = \bigcup^{2^N}_{k=1} I_k
\quad\text{y}\quad
\sum^{2^N}_{k=1} \vol(I_k) = S_N = \left(\frac23\right)^N < \varepsilon
\]
Esto prueba que $\mathcal{C}$ es Jordan-nulo.\vspace{1em}\\
\emph{Veamos que $\mathbb{Q}\cap[0,1]$ es Lebesgue-nulo y no Jordan-nulo}:\\
Enumeremos $\mathbb{Q} \cap [0,1]$: $\set{q_n}_{n\en\mathbb{N}}$, y sea $\varepsilon > 0$ arbitrario. Entonces para cada $n\en\mathbb{N}$ definimos $I_n \coloneq \bigl[q_n-\frac{\varepsilon}{2^{n+2}},q_n+\frac{\varepsilon}{2^{n+2}}\bigr]$. Luego
\[
\mathbb{Q} \cap [0,1] \subseteq \bigcup_{n\en\mathbb{N}} I_n
\quad\text{y}\quad
\sum^\infty_{n=1}\vol(I_n) = \sum^\infty_{n=1}\frac{\varepsilon}{2^{n+1}} = \frac{\varepsilon}{2} < \varepsilon
\]
Esto prueba que $\mathbb{Q} \cap [0,1]$ es Lebesgue-nulo.\par
Por otro lado, para $\epsilon_0 = 1$ se cumple, para cualesquiera $I_1,\dots,I_n$ intervalos compactos con $\mathbb{Q} \cap [0,1] \subseteq \bigcup^n_{k=1} I_k$, que $\sum^n_{k=1} \vol(I_k) \ge 1$. En efecto, como $\bigcup^n_{k=1} I_k$ es cerrado, se tiene 
\[
[0,1] = \ubar{\mathbb{Q} \cap [0,1]} \subseteq \bigcup^n_{k=1} I_k
\]
por lo cual, los $I_k$ son cubierta de intervalos cerrados para $[0,1]$, así que $\sum^n_{k=1} \vol(I_k) \ge 1$ (pues ya sabemos que $\nu([0,1]) = 1$). Esto prueba que $\mathbb{Q} \cap [0,1]$ no es Jordan-nulo.\vspace{1em}\\
\emph{Veamos que $A \subseteq \mathbb{R}^n$ Lebesgue-nulo si y solo si $\nu(A) = 0$}:\\
Por una equivalencia de la def. de ínfimo, $\inf\set[\big]{ \sum^\infty_{k=1} \vol(Q_k) \given Q_k \en \mathcal{Q}_n,\; A \subseteq \bigcup^{\infty}_{k=1} Q_k} = \nu(A) = 0$ si y solo si $0$ es cota inferior del conjunto y
\begin{equation}\label{eq:lebnull}
\forall \varepsilon>0, \exists{Q_1,Q_2,\dots} \en \mathcal{Q}_n,\quad
A \subseteq \bigcup^{\infty}_{k=1} Q_k \quad\text{y}\quad\sum^\infty_{k=1} \vol(Q_k) < \varepsilon
\end{equation}
Como siempre se tiene $0 \leq \sum^\infty_{k=1} \vol(Q_k)$, la condición de que $0$ sea cota inferior del conjunto es redundante, así que lo anterior visto se traduce en que $\nu(A) = 0$ si y solo si se cumple (\ref{eq:lebnull}); en otras palabras, $\nu(A) = 0$ si y solo si $A$ es Lebesgue-nulo.\vspace{1em}\\
\emph{Encontremos un abierto cuya frontera tenga medida positiva}:\\
Consideremos cualquier $\mathcal{C}_\alpha$ con $0 < \alpha < \nicefrac{1}{3}$ y tomemos $U \coloneq \mathbb{R}\setminus\mathcal{C}_\alpha$. Luego, por propiedades de la frontera:
\[
\partial U = \partial \mathcal{C}_\alpha = \ubar{\mathcal{C}_\alpha}\setminus\inter(\mathcal{C}_\alpha)
\]
Pero $\ubar{\mathcal{C}_\alpha}=\mathcal{C}_\alpha$ y es fácil ver que $\inter(\mathcal{C}_\alpha) = \varnothing$, ya que cada $C^\alpha_k$ es unión de intervalos con longitud
\[
2^{-k}S_k
= 2^{-k} - \frac{\alpha^{k+1} - 2^{-k}}{2\alpha-1} = \frac{2^{-k+1}\alpha - \alpha^{k+1}}{2\alpha-1} 
\]
Como esta longitud converge a cero (pues $2^{-k+1} \to 0$ y $\alpha^k \to 0$ ya que $\alpha<1$) se sigue que, dado un intervalo abierto arbitrario $I$, siempre existe un $N$ tal que $2^{-N}S_N$ es menor que la longitud de $I$, por lo que $I$ no puede estar contenido en $C^\alpha_N$ (como $I$ es conexo, tendría que estar contenido en un intervalo de $C_N$, pero todos estos tienen longitud $2^{-N}S_N$, que es menor a la de $I$); como además $\mathcal{C}_\alpha \subseteq C^\alpha_N$, se sigue que también $I \not\subseteq \mathcal{C}_\alpha$. Esto prueba que $\mathcal{C}_\alpha$ no contiene ningún intervalo abierto, así que su interior es vacío.\par
Hemos probado $\partial U = \mathcal{C}_\alpha$. Solo falta ver que $m(\mathcal{C}_\alpha)>0$. En efecto, $\set{C^\alpha_k}_k$ es sucesión decreciente de borelianos (medibles); más aún, la aditividad y el hecho de que la medida de un intervalo es su longitud implica directamente que $m(C^\alpha_k) = S_k < \infty$, por lo que (al ser $m$ medida):
\[
m(\mathcal{C}_\alpha) = m\biggl(\bigcap^\infty_{k=1} C^\alpha_k\biggr) = \lim_{k\to\infty} S_k = \lim_{k\to\infty} 1 - \alpha\frac{(2\alpha)^k - 1}{2\alpha-1}
\]
pero $2\alpha < \nicefrac{2}{3}$, así que $(2\alpha)^k \to 0$, por lo cual
\[
m(\mathcal{C}_\alpha) = 1+ \frac{\alpha}{2\alpha-1} = \frac{3\alpha-1}{2\alpha-1}
\]
como $2\alpha < 3\alpha < 1$, se sigue que $m(\mathcal{C}_\alpha) > 0$.
\end{sol}